\documentclass[a4paper]{article}
% UTF-8
\usepackage[T1]{fontenc}
\usepackage[utf8]{inputenc}

% Czcionka, która sensownie wygląda na monitorze.
\usepackage{lmodern}

% Język
\usepackage[polish]{babel}
% \usepackage[english]{babel}

% Rózne przydatne paczki:
% - znaczki matematyczne
\usepackage{amsmath, amsfonts}
% - wcięcie na początku pierwszego akapitu
\usepackage{indentfirst}
% - komenda \url 
\usepackage{hyperref}
% - dołączanie obrazków
\usepackage{graphics}
% - szersza strona
\usepackage[nofoot,hdivide={2cm,*,2cm},vdivide={2cm,*,2cm}]{geometry}
\frenchspacing
% - brak numerów stron
\pagestyle{empty}

% dane autora
\author{Mateusz Nawrot}
\title{Sprawozdanie z listy 7}
\date{\today}

% początek dokumentu
\begin{document}
\maketitle
\section{Czas uniksowy i problem roku 2038}
Czas uniksowy liczy sekundy od północy 1 stycznia 1970 roku w strefie UTC.
Niektóre programy zapisują tę liczbę w 32 bitach ze znakiem.
19 stycznia 2038 roku liczba sekund przekroczy zakres tego zapisu, więc program
może odczytać błędną datę. Rozwiązaniem jest zapis czasu w 64 bitach.
\section{Dowód indukcyjny}
Chcemy wykazać, że dla każdej liczby całkowitej \(n\geq 0\) zachodzi
\[
\sum_{i=0}^{n}2^i+1=2^{n+1}.
\]
\textbf{Baza indukcji.} Dla \(n=0\) lewa strona jest równa \(2^0+1=2\), a prawa \(2^1=2\).

\textbf{Założenie indukcyjne.} Przyjmijmy, że dla pewnego \(k\geq 0\) całkowitego zachodzi \(\sum_{i=0}^k 2^i+1=2^{k+1}\).

\textbf{Krok indukcyjny.} W sumie dla \(n=k+1\) wydzielamy ostatni składnik i korzystamy z założenia indukcyjnego:
\[
\sum_{i=0}^{k+1}2^i+1
=\left(\sum_{i=0}^{k}2^i+1\right)+2^{k+1}
=2^{k+1}+2^{k+1}=2^{k+2}.
\]
Ponieważ \(k+2=(k+1)+1\), otrzymujemy tezę dla \(k+1\).

\textbf{Wniosek.} Na mocy zasady indukcji matematycznej wzór jest prawdziwy dla każdej liczby całkowitej \(n\geq 0\).
\section{Pies}
\begin{verbatim}
       / \__
      (    @\___
      /         O
     /   (_____/
    /_____/   U
\end{verbatim}
\end{document}
